% ---------------------------------------------------------------------------
%  QML Lab research project 16: Simulator vs IBM Hardware Benchmark
%  Build:  latexmk -pdf main.tex   (needs qmlproject.sty, harvardqml.bst, references.bib)
% ---------------------------------------------------------------------------
\documentclass[11pt,a4paper]{article}
\usepackage{qmlproject}
\qmlheader{Project 16: Simulator vs hardware}

\hypersetup{pdftitle={QML Lab project 16: Simulator vs IBM Hardware Benchmark}}

\begin{document}

\begin{center}
  {\small\textcolor{black!55}{QML Lab \textbullet\ research project 16}}\\[0.5em]
  {\LARGE\bfseries Simulator vs IBM Hardware Benchmark}\\[0.6em]
  {\large Measuring and explaining the gap between simulated and real-device results.}\\[0.6em]
  {\small Living document \textbullet\ last updated \today}
\end{center}

\begin{projectcard}
\small
\begin{tabularx}{\linewidth}{@{}>{\bfseries}l X >{\bfseries}l X@{}}
Quadrant  & \CQ\ (classical data, quantum processor) & Track & Noise and hardware \\
Difficulty & Medium & Tools & Qiskit, Qiskit Runtime, Qiskit Aer \\
Builds on & Modules 2 and 8 of the QML Lab & Workspace & \path{qml_projects/simulator-vs-ibm-hardware/} \\
\end{tabularx}
\end{projectcard}

\paragraph{Summary.}
Every result in the lab so far comes from simulators. This project runs the best small model on a real
IBM quantum processor, measures the gap to the simulator, and builds an error budget that explains where
the gap comes from: readout errors, gate errors, transpilation and drift. It produces the lab's reference
for what the other projects can expect on hardware.

% ===========================================================================
\section{Background and key concepts}
% ===========================================================================
\subsection{From circuit to device}
A circuit must be transpiled to the device's native gates and qubit connectivity, which adds SWAP gates
and depth. Qiskit \citep{javadiabhari2024} handles this, and Qiskit Runtime executes circuits through
primitives (Sampler and Estimator) with optional error suppression and mitigation. Current devices are
noisy intermediate-scale processors without full error correction \citep{preskill2018}.

\subsection{Error sources}
\begin{description}[style=nextline,leftmargin=1.2em]
  \item[Gate errors] average error rates per gate, measured by randomised benchmarking \citep{magesan2011};
        two-qubit gates dominate.
  \item[Readout errors] described by a confusion matrix $A_{ij}=P(\text{read }i\,|\,\text{prepared }j)$; observed
        distributions are $p_{\text{obs}}=Ap_{\text{true}}$ and can be partly inverted \citep{bravyi2021readout}.
  \item[Decoherence] energy relaxation ($T_1$) and dephasing ($T_2$) limit circuit duration.
  \item[Drift and crosstalk] calibration changes over hours and days; neighbouring qubits interfere.
\end{description}

\subsection{Error budget}
The accuracy gap $\Delta=\text{acc}_{\text{sim}}-\text{acc}_{\text{hw}}$ is decomposed by simulating the
same circuits with noise models that switch on one error source at a time, using the device's own
calibration data. Hardware-efficient circuits \citep{kandala2017} keep the transpiled depth low.
\citet{kim2023} showed that with careful mitigation, current processors can produce accurate expectation
values on circuits beyond brute-force simulation, which sets the state of the art for what is achievable.

% ===========================================================================
\section{Research questions}
% ===========================================================================
The project is open-ended: these questions are the starting point, not the finish line. Add new
ones to the README and the journal as they come up.

\begin{rqlist}
  \item How large is the accuracy gap between simulator and hardware, and how much of it is readout versus gate error?
  \item How much does transpilation (layout, optimisation level) change the gap?
  \item How stable are hardware results from one day to the next?
\end{rqlist}

% ===========================================================================
\section{Method}
% ===========================================================================
\subsection{Data}
The test set of the chosen model (Module 5 or Project 04 classifier with 2--4 qubits, 100 test points);
trained weights are fixed beforehand on the simulator. Each run records backend name, date, calibration
snapshot, transpiled depth and two-qubit gate count.

\subsection{Models}
Evaluation pipeline in Qiskit: exact state vector, shot-based simulation, Aer with noise models built from
backend calibration (readout only, gates only, full), and the real backend via Runtime with and without
readout mitigation. IBM Quantum credentials are read from the local Qiskit account, never stored in the
repository.

\subsection{Baselines}
Exact simulation (upper bound); shot-based simulation with the same number of shots (pure sampling
error); Aer with the full device noise model (prediction of hardware performance).

\subsection{Metrics}
\begin{itemize}
  \item Test accuracy and expectation-value error against exact values.
  \item Accuracy gap $\Delta$ and its decomposition by error source.
  \item Transpiled depth, two-qubit gate count and QPU time per run.
\end{itemize}

% ===========================================================================
\section{Planned experiments}
% ===========================================================================
Each experiment becomes one \texttt{[[experiments]]} entry in the workspace, so its runs, parameters
and metrics are stored and compared in the portal.

\begin{table}[h]
\centering\small
\begin{tabularx}{\linewidth}{@{}l>{\raggedright}X>{\raggedright}X>{\raggedright\arraybackslash}X@{}}
\toprule
ID & Experiment & What varies & Main metrics \\
\midrule
  E1 & Sampling only & exact vs shots & accuracy, error \\
  E2 & Noise components & readout / gates / full model & predicted $\Delta$ \\
  E3 & Hardware run & real backend & measured $\Delta$ \\
  E4 & Mitigation & resilience options & $\Delta$, overhead \\
  E5 & Drift & same job on 3 days & spread of results \\
\bottomrule
\end{tabularx}
\caption{Experiment plan. Results go into the journal and the results log (\cref{sec:log}).}
\end{table}

% ===========================================================================
\section{Running it in the QML Lab}
% ===========================================================================
The workspace lives in \path{qml_projects/simulator-vs-ibm-hardware/}. Scripts import three helpers from
\path{qml_lab/plots.py}:
\begin{itemize}
  \item \texttt{params()} reads the parameters from the portal's run form (defaults when run by hand);
  \item \texttt{save()} stores figures with the run;
  \item \texttt{log\_metrics()} records numbers that appear in the run table.
\end{itemize}

\begin{lstlisting}
qml_projects/simulator-vs-ibm-hardware/
  project.toml          metadata + experiments
  README.md             goal, open questions, ideas
  experiments/          runnable scripts
  docs/main.tex         this document
\end{lstlisting}

First experiment: This experiment is \emph{planned}: write the script, add the block below to \texttt{project.toml}, then run \texttt{python manage.py sync\_projects}.

\begin{lstlisting}[language={}]
[[experiments]]
slug = "hardware-gap"
title = "Simulator vs hardware accuracy for the Module 5 classifier"
script = "experiments/e01_hardware_gap.py"
figure = "gap_breakdown.png"

[experiments.params]
backend = "least_busy"
shots = 4000
optimization_level = 3
mitigate_readout = true
n_test = 100
seed = 0
\end{lstlisting}

\begin{lstlisting}[language=bash]
# from the portal folder: pick up changes to project.toml
python manage.py sync_projects
cd qml_projects/simulator-vs-ibm-hardware
# run by hand with the default parameters
PYTHONPATH=../../qml_lab python experiments/e01_hardware_gap.py
\end{lstlisting}

% ===========================================================================
\section{Quality bar}
% ===========================================================================
\begin{itemize}
  \item A classical baseline with a similar number of trainable parameters, tuned with the same care \citep{bowles2024}.
  \item At least five random seeds per configuration, reported as mean $\pm$ standard deviation.
  \item How results change with the number of qubits.
  \item At least one run under a realistic noise model.
  \item Record backend, date and calibration with every hardware result; without them a number cannot be reproduced.
\end{itemize}

% ===========================================================================
\section{Open directions}
% ===========================================================================
\begin{itemize}
  \item Compare two backends or processor generations.
  \item Run a quantum kernel (Project 01) and a trained policy (Project 10) on hardware.
  \item Add zero-noise extrapolation from Project 14.
\end{itemize}

% ===========================================================================
\section{Risks and limitations}
% ===========================================================================
\begin{itemize}
  \item Queue times and free-plan limits restrict the number of runs; plan circuits carefully.
  \item Calibration drift makes single runs unreliable; repeat across days.
  \item Never commit API tokens; use the local Qiskit account configuration.
\end{itemize}

% ===========================================================================
\section{Results log}
\label{sec:log}
% ===========================================================================
Summarise findings here as they become stable; day-to-day notes belong in the portal journal.

\begin{table}[h]
\centering\small
\begin{tabularx}{\linewidth}{@{}l l >{\raggedright}X >{\raggedright\arraybackslash}X@{}}
\toprule
Date & Run & Finding & Next step \\
\midrule
 & & & \\
 & & & \\
 & & & \\
\bottomrule
\end{tabularx}
\end{table}

\bibliographystyle{harvardqml}
\bibliography{references}

\end{document}
